% SPDX-License-Identifier: CC0-1.0
\documentclass[11pt,a4paper]{article}

\usepackage[T1]{fontenc}
\usepackage[utf8]{inputenc}
\usepackage{mathpazo}
\usepackage{microtype}
\usepackage[a4paper,margin=27mm,headheight=23pt]{geometry}
\usepackage{amsmath,amssymb,amsthm,mathtools}
\usepackage{booktabs,tabularx,array,longtable}
\usepackage{enumitem}
\usepackage{xcolor}
\usepackage{fancyhdr}
\usepackage{lastpage}
\usepackage{hyperref}
\usepackage{xurl}
\usepackage[nameinlink,noabbrev]{cleveref}

\definecolor{ink}{HTML}{172033}
\definecolor{accent}{HTML}{245A86}
\definecolor{accentlight}{HTML}{EAF2F8}
\definecolor{warning}{HTML}{8A4B08}
\definecolor{warninglight}{HTML}{FFF5E6}
\definecolor{rulegray}{HTML}{CBD5E1}

\hypersetup{
  colorlinks=true,
  linkcolor=accent,
  citecolor=accent,
  urlcolor=accent,
  pdfauthor={GPT 5.6 Sol; Claude Opus 4.8 / 5},
  pdftitle={Two vertex-disjoint monochromatic K4s in every two-colouring of K20},
  pdfsubject={A candidate certificate-backed determination of VR2(K4)},
  pdfkeywords={Ramsey theory, vertex-disjoint Ramsey number, K4, SAT certificate, DRUP, LRAT}
}
\urlstyle{same}

\pagestyle{fancy}
\fancyhf{}
\lhead{\small\textcolor{ink}{Two vertex-disjoint monochromatic \(K_4\)s in \(K_{20}\)}}
\rhead{\small\textcolor{ink}{Candidate certificate-backed result}}
\cfoot{\small\textcolor{ink}{\thepage\ of \pageref*{LastPage}}}
\renewcommand{\headrulewidth}{0.4pt}
\renewcommand{\headrule}{\hbox to\headwidth{\color{rulegray}\leaders\hrule height \headrulewidth\hfill}}

\setlength{\parindent}{0pt}
\setlength{\parskip}{0.58em}
\setlist[itemize]{leftmargin=1.35em,itemsep=0.22em,topsep=0.35em}
\setlist[enumerate]{leftmargin=1.6em,itemsep=0.28em,topsep=0.35em}
\renewcommand{\arraystretch}{1.16}

\newtheorem{theorem}{Theorem}[section]
\newtheorem{lemma}[theorem]{Lemma}
\newtheorem{proposition}[theorem]{Proposition}
\newtheorem{corollary}[theorem]{Corollary}
\theoremstyle{definition}
\newtheorem{definition}[theorem]{Definition}
\theoremstyle{remark}
\newtheorem{remark}[theorem]{Remark}

\newcommand{\homfour}{\textnormal{hom-}4}
\newcommand{\corezero}{\texttt{\detokenize{OsHHirKdlp[IFVI|KpqfR}}}
\newcommand{\coreone}{\texttt{\detokenize{Ov?IXZIhlRWjUXL[iphst}}}
\newcommand{\sha}[1]{\texttt{\footnotesize #1}}
\newcommand{\artifact}[1]{\texttt{\detokenize{#1}}}
\newcommand{\statusbox}[1]{%
  \begin{center}
  \fcolorbox{warning}{warninglight}{%
    \parbox{0.91\textwidth}{\small\textcolor{warning}{\textbf{Status.}} #1}}
  \end{center}}
\newcommand{\resultbox}[1]{%
  \begin{center}
  \fcolorbox{accent}{accentlight}{%
    \parbox{0.90\textwidth}{\centering\large\textcolor{accent}{#1}}}
  \end{center}}

\title{\vspace{-1.1cm}\textbf{\color{ink}Two vertex-disjoint monochromatic \(K_4\)s\\
 in every two-colouring of \(K_{20}\)}\\[0.42em]
\large\color{accent}A candidate certificate-backed determination of \(\mathrm{VR}_2(K_4)\)}
\author{GPT 5.6 Sol \and Claude Opus 4.8 / 5}
\date{26 July 2026}

\begin{document}
\maketitle
\thispagestyle{plain}

\begin{abstract}
For a graph \(F\) without isolated vertices and a positive integer \(t\),
let \(\mathrm{VR}_t(F)\)
denote the least \(n\) such that every red--blue colouring of \(K_n\)
contains \(t\) pairwise vertex-disjoint monochromatic copies of \(F\);
the copies need not have the same colour.  We give a candidate
computer-assisted proof that
\[
  \mathrm{VR}_2(K_4)=20.
\]
The lower bound is hand-checkable.  It is supplied by a \(19\)-vertex
colouring obtained from the Paley graph \(P_{17}\) by replacing one vertex
with three independent twins.  For the upper bound, fix one homogeneous
\(4\)-set in a hypothetical counterexample on \(20\) vertices.  The graph
induced by the remaining \(16\) vertices must be a \((4,4)\)-Ramsey graph,
of which there are exactly two up to isomorphism.  Each fixed core leaves \(64\)
undetermined cross edges.  Enumerating the \(872\) mixed four-sets that
can be homogeneous and forbidding every disjoint pair gives a
\(64\)-variable CNF with \(104{,}524\) clauses.  Both formulas have
RUP/DRAT/LRAT-checked refutations.  The paper separates the graph-theoretic
reduction, the exact encoding equivalence, the proof-object checks, and
the remaining trusted-computing assumptions.
\end{abstract}

\noindent\textbf{Keywords:} Ramsey theory; vertex-disjoint Ramsey number;
monochromatic clique; exhaustive generation; SAT certificate; DRUP; LRAT.

\noindent\textbf{2020 Mathematics Subject Classification:} 05C55, 05C85.

\statusbox{This is an \textbf{unrefereed candidate computer-assisted proof}
offered for independent mathematical and computational scrutiny.  The
lower-bound construction is hand-verifiable.  The two upper-bound proof
objects have been replayed locally by separately implemented checkers, including a
formally verified CakeML checking core.  No external reproduction, human
peer review, or end-to-end proof-assistant formalisation is claimed.}

\section{Introduction}

For a graph \(F\) without isolated vertices and a positive integer \(t\),
the vertex-disjoint Ramsey number \(\mathrm{VR}_t(F)\) is the least order
of a complete graph whose every red--blue edge-colouring contains \(t\)
pairwise vertex-disjoint monochromatic copies of \(F\)
\cite{jent-osborn-zhang,jent-zhang}.  Distinct copies may receive different
colours.  This distinguishes \(\mathrm{VR}_t(F)\) from the ordinary
symmetric Ramsey number of the disconnected graph \(tF\), which requires
a monochromatic union of all \(t\) copies.

Recent work has determined such numbers for triangles, paths, stars,
matchings, and other graphs of small order or size
\cite{jent-osborn-zhang,jent-zhang}.  We consider \(K_4\).

\begin{theorem}[Main result]\label{thm:main}
Every red--blue colouring of \(K_{20}\) contains two vertex-disjoint
monochromatic copies of \(K_4\), while some colouring of \(K_{19}\) does
not.  Equivalently,
\[
  \boxed{\mathrm{VR}_2(K_4)=20}.
\]
\end{theorem}

The proof has two unequal parts.  The lower bound is a direct construction
based on the Paley graph \(P_{17}\).  The upper bound uses the classical
identity \(R(4,4)=18\) and a complete catalogue of the two
\((4,4)\)-graphs on \(16\) vertices.  Fixing either graph leaves a
\(64\)-variable propositional problem, and both resulting formulas are
unsatisfiable.

The upper-bound computation originally appeared as the structural lemma
inside a candidate proof concerning the cochromatic number of graphs on
\(20\) vertices.  The present paper extracts that lemma, supplies the
sharp \(19\)-vertex construction, and states the result in the established
vertex-disjoint Ramsey notation.  The cocolouring computations in the
larger proof are not used here.

The logical architecture is as follows.

\begin{center}
\begin{tabularx}{0.96\textwidth}{@{}>{\bfseries}p{0.23\textwidth} X@{}}
\toprule
Component & Role\\
\midrule
Lower construction &
A hand-checkable colouring of \(K_{19}\) in which every monochromatic
\(K_4\) uses at least two vertices of a fixed three-vertex set.\\
Ramsey reduction &
A counterexample on \(20\) vertices would consist of one homogeneous
\(4\)-set, one of the two \(16\)-vertex \((4,4)\)-cores, and \(64\)
cross edges.\\
Exact encoding &
The CNF is satisfied exactly by cross-edge assignments with no two
disjoint homogeneous \(4\)-sets.\\
Certificates &
Separately implemented local checkers have accepted both fixed-core
refutations.\\
\bottomrule
\end{tabularx}
\end{center}

\section{Definitions and elementary invariance}

Represent a red--blue colouring of \(K_n\) by a finite simple graph \(G\):
red edges are the edges of \(G\), and blue edges are the edges of
\(\overline G\).

\begin{definition}
A vertex set \(S\subseteq V(G)\) is \emph{homogeneous} if \(G[S]\) is
complete or edgeless.  Thus a homogeneous \(4\)-set is exactly a
monochromatic \(K_4\) in the corresponding two-colouring.  A
\((4,4)\)-graph is a graph containing neither a \(K_4\) nor an independent
set of order four.
\end{definition}

\begin{lemma}[Complement invariance]\label{lem:complement}
A vertex set is homogeneous in \(G\) if and only if it is homogeneous in
\(\overline G\).  Consequently, the property of containing two
vertex-disjoint homogeneous \(4\)-sets is invariant under complementation.
\end{lemma}

\begin{proof}
Complementation exchanges cliques and independent sets without changing
their vertex sets.
\end{proof}

\section{A sharp lower-bound colouring on nineteen vertices}

Let
\[
  Q=\{1,2,4,8,9,13,15,16\}\subset\mathbb Z_{17}
\]
be the set of nonzero quadratic residues.  The Paley graph \(P_{17}\) has
vertex set \(\mathbb Z_{17}\), with \(i\sim j\) when
\(i-j\pmod {17}\in Q\) \cite{paley}.  Since \(-1=16\in Q\), the graph is
undirected.

\begin{lemma}\label{lem:paley}
The graph \(P_{17}\) contains neither a \(K_4\) nor an independent set of
order four.
\end{lemma}

\begin{proof}
Translations and multiplications by nonzero squares are automorphisms of
\(P_{17}\), and they act transitively on its edges.  A hypothetical
\(K_4\) may therefore be assumed to contain \(\{0,1\}\).  The common
neighbours of \(0\) and \(1\) are
\[
  Q\cap(1+Q)=\{2,9,16\}.
\]
Their pairwise differences are congruent, up to sign, to \(7\) or \(14\),
neither of which lies in \(Q\).  They are therefore pairwise nonadjacent,
so no two can complete \(\{0,1\}\) to a \(K_4\).

Multiplication by a quadratic nonresidue, for example \(3\), is an
isomorphism \(P_{17}\cong\overline{P_{17}}\).  The absence of a \(K_4\)
therefore also gives the absence of an independent four-set.
\end{proof}

We now replace vertex \(9\) by three independent twins.  More precisely,
start with \(P_{17}-9\), add a set
\[
  X=\{x_1,x_2,x_3\},
\]
make \(X\) independent, and give every \(x_i\) the former external
neighbourhood of vertex \(9\):
\[
  N=\{0,1,5,7,8,10,11,13\}.
\]
The remaining outside vertices are
\[
  M=\{2,3,4,6,12,14,15,16\}.
\]
Call the resulting graph \(L_{19}=P_{17}[9\rightsquigarrow\overline K_3]\).

\begin{lemma}\label{lem:local-paley}
The graph \(P_{17}[N]\) is triangle-free, and \(P_{17}[M]\) has no
independent triple.
\end{lemma}

\begin{proof}
The edges of \(P_{17}[N]\) form the eight-cycle
\[
  0-8-7-5-1-10-11-13-0
\]
together with the four opposite chords
\[
  01,\qquad 8\,10,\qquad 7\,11,\qquad 5\,13.
\]
This is the eight-vertex M\"obius ladder, hence is triangle-free; the
claim can also be checked directly from the displayed edge list.

The edges of \(\overline{P_{17}[M]}\) form the eight-cycle
\[
  2-12-15-4-16-6-3-14-2
\]
together with
\[
  2\,16,\qquad 12\,6,\qquad 15\,3,\qquad 4\,14.
\]
This graph is also triangle-free.  Therefore \(P_{17}[M]\) has no
independent triple.
\end{proof}

\begin{proposition}\label{prop:lower}
The colouring represented by \(L_{19}\) contains no two
vertex-disjoint monochromatic \(K_4\)s.  Hence
\[
  \mathrm{VR}_2(K_4)\geq20.
\]
\end{proposition}

\begin{proof}
There is no red \(K_4\).  A red \(K_4\) wholly outside \(X\) would lie in
\(P_{17}-9\), contradicting Lemma~\ref{lem:paley}.  A red \(K_4\) using one
vertex of \(X\) would require a triangle in \(N\), contrary to
Lemma~\ref{lem:local-paley}.  A red \(K_4\) using at least two vertices of
\(X\) is impossible because \(X\) is independent.

A blue \(K_4\) wholly outside \(X\) would be an independent four-set in
\(P_{17}-9\).  A blue \(K_4\) using exactly one vertex of \(X\) would
require an independent triple in \(M\).  Both possibilities are excluded
by Lemmas~\ref{lem:paley} and~\ref{lem:local-paley}.  Consequently every
blue \(K_4\)
contains at least two of the three vertices of \(X\).  Any two such
four-sets intersect.  Thus the colouring contains no two vertex-disjoint
monochromatic \(K_4\)s.
\end{proof}

\begin{remark}[Sanity check]\label{rem:lower-count}
The graph \(L_{19}\) has \(84\) red edges and degree sequence
\(8^{11}10^8\).  It contains no red \(K_4\) and exactly \(44\) blue
\(K_4\)s.  Indeed, a blue copy uses either two twins and one of the
\(12\) nonedges of \(P_{17}[M]\), or all three twins and one vertex of
\(M\):
\[
  \binom32\!\cdot12+\binom33\!\cdot8=44.
\]
A direct verifier is included with the source bundle.
\end{remark}

\section{Reduction of the twenty-vertex upper bound}

The classical Ramsey identity \(R(4,4)=18\) was established by Greenwood
and Gleason \cite{greenwood-gleason}.  It is also reproduced by the
catalogue computation in \cref{sec:catalogue}.

\begin{lemma}\label{lem:two-homfour}
Every graph on \(20\) vertices contains two vertex-disjoint homogeneous
\(4\)-sets.
\end{lemma}

\begin{proof}
Suppose that a graph \(G\) on \(20\) vertices contains no two disjoint
homogeneous \(4\)-sets.  Since \(20\geq R(4,4)=18\), the graph contains a
homogeneous \(4\)-set \(A\).  By Lemma~\ref{lem:complement}, complementing
\(G\) if necessary preserves the target property, so assume that \(A\)
is a clique.  Relabel
\[
  A=\{a_0,a_1,a_2,a_3\},
  \qquad B=V(G)\setminus A.
\]
If \(G[B]\) contains a homogeneous \(4\)-set, it is disjoint from \(A\)
and the conclusion follows.  Thus a counterexample must have \(G[B]\) a
\((4,4)\)-graph on \(16\) vertices.

The complete catalogue in \cref{sec:catalogue} contains exactly two such
graphs up to isomorphism.  In graph6 notation they are
\[
  C_0=\corezero,
  \qquad
  C_1=\coreone.
\]
Each has \(60\) edges, degree sequence \(7^8 8^8\), \(56\) triangles,
and \(56\) independent triples.

Once \(G[B]\) is fixed to \(C_0\) or \(C_1\), the only undetermined edges
are the \(4\cdot16=64\) edges between \(A\) and \(B\).  For each core,
\cref{sec:encoding} constructs a CNF whose satisfying assignments are
exactly those cross-edge patterns for which the resulting graph contains
no two disjoint homogeneous \(4\)-sets.  Both formulas are unsatisfiable
by the certificate checks in \cref{sec:certificates}.  This contradiction
proves the lemma.
\end{proof}

\section{The exact fixed-core encoding}\label{sec:encoding}

Fix \(A=\{a_0,a_1,a_2,a_3\}\) as a clique and let
\(B=\{b_0,\ldots,b_{15}\}\) induce one of the two cores.  For
\(0\leq i<4\) and \(0\leq j<16\), introduce a Boolean variable
\[
  x_{i,j}=1\quad\Longleftrightarrow\quad a_i b_j\in E(G).
\]

\begin{lemma}[Mixed-set classification]\label{lem:mixed}
The mixed four-sets that can be homogeneous are exactly the four types in
\cref{tab:mixed}.
\end{lemma}

\begin{table}[ht]
\centering
\caption{Conditions under which a mixed four-set is homogeneous.}
\label{tab:mixed}
\begin{tabularx}{\textwidth}{@{}l X r@{}}
\toprule
Split & Necessary and sufficient condition & Count\\
\midrule
\(3A+1B\) &
Clique: the three relevant cross edges are present. &
\(\binom43\cdot16=64\)\\
\(2A+2B\) &
Clique: the \(B\)-pair is one of the \(60\) core edges and all four
relevant cross edges are present. &
\(\binom42\cdot60=360\)\\
\(1A+3B\) &
Clique: the \(B\)-triple is one of the \(56\) triangles and all three
cross edges are present. &
\(4\cdot56=224\)\\
\(1A+3B\) &
Independent: the \(B\)-triple is one of the \(56\) independent triples
and all three cross edges are absent. &
\(4\cdot56=224\)\\
\midrule
&\textbf{Total}&\(\mathbf{872}\)\\
\bottomrule
\end{tabularx}
\end{table}

\begin{proof}
A four-set containing three vertices of \(A\) already contains a triangle,
so it can only be a clique; this occurs exactly when its three cross edges
are present.  A four-set containing two vertices of \(A\) contains their
edge and therefore cannot be independent.  It is a clique exactly when
the \(B\)-pair is an edge and all four cross edges are present.

With one vertex in \(A\), the \(B\)-triple must be either a triangle,
followed by three present cross edges, or an independent triple, followed
by three absent cross edges.  No four-set wholly in \(B\) is homogeneous,
and \(A\) is the only four-set wholly in \(A\).
\end{proof}

Represent each possible mixed homogeneous set \(S\) by the conjunction
\(C_S\) of its three or four required literals.  Whenever
\(S\cap T=\varnothing\), add the clause
\[
  \neg C_S\lor\neg C_T.
\]
For example, if
\[
  S=\{a_0,a_1,a_2,b_5\},
  \qquad
  T=\{a_3,b_1,b_4,b_9\},
\]
and \(\{b_1,b_4,b_9\}\) is an independent triple of the core, the clause
is
\[
 \neg x_{0,5}\vee\neg x_{1,5}\vee\neg x_{2,5}\vee
 x_{3,1}\vee x_{3,4}\vee x_{3,9}.
\]

\begin{proposition}[Encoding equivalence]\label{prop:encoding}
For either core, an assignment to the \(64\) variables satisfies the
constructed CNF if and only if the corresponding \(20\)-vertex graph has
no two vertex-disjoint homogeneous \(4\)-sets.
\end{proposition}

\begin{proof}
By Lemma~\ref{lem:mixed}, every mixed set that can be homogeneous is represented
by exactly one conjunction \(C_S\), and \(C_S\) holds exactly when \(S\)
is homogeneous.  The clause belonging to disjoint sets \(S,T\) is false
exactly when both sets are homogeneous.  Conjoining the clauses over all
disjoint pairs therefore forbids exactly, and only, the existence of two
disjoint mixed homogeneous four-sets.

In the branch being encoded, \(B\) contains no homogeneous four-set, so
\(A\) cannot be paired with a disjoint homogeneous set.  Every prohibited
pair is therefore covered by the mixed-set clauses.
\end{proof}

For each core, the resulting DIMACS formula has \(64\) variables and
\(104{,}524\) clauses.  Its clause-length histogram is
\[
  43{,}840\text{ of length }6,\qquad
  52{,}416\text{ of length }7,\qquad
  8{,}268\text{ of length }8.
\]

\section{Catalogue generation}\label{sec:catalogue}

McKay's maintained Ramsey catalogue lists two non-isomorphic
\((4,4)\)-graphs of order \(16\) \cite{mckay-catalogue}.  The proof bundle
also regenerates the catalogue from the definition, using nauty for
canonical labelling \cite{mckay-piperno}.

\subsection{Completeness of the generation scheme}

Suppose a complete set of isomorphism-class representatives is known at
order \(n\).  Every \((4,4)\)-graph \(H\) on \(n+1\) vertices remains a
\((4,4)\)-graph after deleting a vertex.  Therefore \(H\) arises from an
order-\(n\) representative by adjoining one vertex with some neighbourhood
\(S\).  Trying every admissible \(S\), canonicalising each extension, and
retaining one copy of every canonical form reaches every order-\(n+1\)
class.  Induction begins with the unique one-vertex graph.

For a new vertex with neighbourhood \(S\), a new homogeneous four-set must
contain that vertex.  The extension is admissible exactly when
\[
  G[S]\text{ is triangle-free}
  \quad\text{and}\quad
  G[V\setminus S]\text{ has no independent triple}.
\]
The implementation restricts to \(n-8\leq |S|\leq8\).  This pruning follows
from the elementary bound that a triangle-free graph with independence
number at most three has at most eight vertices.  One proof uses
\(R(3,3)\leq6\): in a hypothetical nine-vertex example every degree is at
most three; a vertex of degree at most two would leave six non-neighbours,
which contain an independent triple, while all degrees equal to three is
impossible by the handshaking lemma.  Applying the same argument to the
complement gives the lower bound on \(|S|\).

\subsection{Replayed output}

The complete run was replayed on 26 July 2026.  The numbers of
non-isomorphic \((4,4)\)-graphs were
\[
\begin{array}{c|rrrrrrrrr}
 n&1&2&3&4&5&6&7&8&9\\ \hline
 \#&1&2&4&9&24&84&362&2{,}079&14{,}701
\end{array}
\]
\[
\begin{array}{c|rrrrrrrrr}
 n&10&11&12&13&14&15&16&17&18\\ \hline
 \#&103{,}706&546{,}356&1{,}449{,}166&1{,}184{,}231&
 130{,}816&640&2&1&0.
\end{array}
\]
The two order-\(16\) outputs agree character-for-character with the
catalogue strings \(C_0,C_1\).  The single survivor at order \(17\) and
empty level at order \(18\) locally reproduce \(R(4,4)=18\).

\section{Certificate verification}\label{sec:certificates}

For each fixed core, a clause-learning-free DPLL program produced a
tree-like DRUP refutation.  Core \(0\) visited \(1{,}571\) search nodes,
made \(785\) decisions, and emitted \(1{,}571\) proof lines.  Core \(1\)
visited \(789\) nodes, made \(394\) decisions, and emitted \(789\) proof
lines.  The preserved formulas have SHA-256 hashes
\[
\begin{array}{ll}
C_0:&
\sha{cbfea7b0b2cb712ad8b4f8b129f74c7f010e879064414eaef76e74aa88be3f44},\\[0.28em]
C_1:&
\sha{d3d96a3c2a719c753242673b81aabc3977c5e435bb5996b4c58c85a82d11dc5a}.
\end{array}
\]

The following checks were replayed against formulas regenerated
separately from the two graph6 strings.

\begin{longtable}{@{}>{\raggedright\arraybackslash}p{0.29\textwidth}
>{\raggedright\arraybackslash}p{0.30\textwidth}
>{\raggedright\arraybackslash}p{0.30\textwidth}@{}}
\caption{Replayed validation of the fixed-core refutations.}
\label{tab:verification}\\
\toprule
Check & Core \(0\) & Core \(1\)\\
\midrule
\endfirsthead
\toprule
Check & Core \(0\) & Core \(1\)\\
\midrule
\endhead
Separate formula construction &
\(872\) conditions; \(104{,}524\) clauses; clause set identical &
\(872\) conditions; \(104{,}524\) clauses; clause set identical\\
Direct semantic spot-check &
\(300/300\) assignments agreed &
\(300/300\) assignments agreed\\
Separate RUP checker &
\(1{,}571\) steps; empty clause derived &
\(789\) steps; empty clause derived\\
Package RUP checker &
\texttt{VERIFIED} &
\texttt{VERIFIED}\\
\texttt{drat-trim} \cite{drat-trim} &
\texttt{s VERIFIED}; \(0\) RAT lemmas &
\texttt{s VERIFIED}; \(0\) RAT lemmas\\
\texttt{cake\_lpr} \cite{cake-lpr} &
\texttt{s VERIFIED UNSAT} &
\texttt{s VERIFIED UNSAT}\\
\bottomrule
\end{longtable}

The separately regenerated formulas were byte-identical to the packaged
formulas.  The package audit also tested \(1{,}000\) deterministic
assignments per core, and the twelve-entry SHA-256 manifest passed in
full.  The DRUP proofs were converted to LRAT with
\texttt{drat-trim -L}; the resulting LRAT proofs were accepted against the
original CNFs by \texttt{cake\_lpr}, whose checking core is formally
verified in CakeML \cite{cake-lpr}.

Cross-pairing the first formula with the second proof was rejected by the
standalone RUP checker, by \texttt{drat-trim}, and by \texttt{cake\_lpr}.
One operational detail matters: \texttt{cake\_lpr} can print a checking
failure while returning process status zero, so automated replay must
require the exact verdict \texttt{s VERIFIED UNSAT}.

\section{Completion of the proof and a higher-copy consequence}

\begin{proof}[Proof of \cref{thm:main}]
The construction in Proposition~\ref{prop:lower} gives
\(\mathrm{VR}_2(K_4)\geq20\).  By Lemma~\ref{lem:two-homfour}, every colouring of \(K_{20}\) has
two vertex-disjoint monochromatic \(K_4\)s.  Hence
\(\mathrm{VR}_2(K_4)\leq20\).
\end{proof}

\resultbox{\(\displaystyle \boxed{\mathrm{VR}_2(K_4)=20}\)}

The exact two-copy value immediately improves the standard greedy bound
for every larger number of copies.

\begin{corollary}\label{cor:higher}
For every integer \(t\geq2\),
\[
  \mathrm{VR}_t(K_4)\leq4t+12.
\]
\end{corollary}

\begin{proof}
The case \(t=2\) is Theorem~\ref{thm:main}.  For \(t\geq3\), consider a
colouring of \(K_{4t+12}\).  Since \(4t+12\geq R(4,4)=18\), choose and
remove one monochromatic \(K_4\).  The remaining graph has
order \(4(t-1)+12\), so induction supplies \(t-1\) further pairwise
disjoint monochromatic copies.  This is the general extension principle
recorded in \cite{jent-osborn-zhang}.
\end{proof}

In particular, \(\mathrm{VR}_3(K_4)\leq24\).  A separate classification of
the \(19\)-vertex two-copy extremals sharpens the three-copy problem
further, but that classification and its additional certificates are not
needed for the present theorem.

\section{Assurance boundary and reproducibility}\label{sec:assurance}

The certificate checks establish unsatisfiability of two exact CNFs.  The
full graph-theoretic theorem also uses the catalogue and the semantic
bridge from graphs to those formulas.  The layers should therefore be
distinguished.

\begin{center}
\begin{tabularx}{\textwidth}{@{}>{\bfseries}p{0.22\textwidth} X@{}}
\toprule
Layer & What is established, and what remains trusted\\
\midrule
Hand mathematics &
The Paley construction, lower-bound proof, complement invariance,
catalogue-completeness induction, mixed-set classification, encoding
equivalence, and final deductions can be checked directly from the text.\\
SAT proof objects &
The two DRUP refutations are accepted by separately implemented local RUP checkers
and by \texttt{drat-trim}; their LRAT translations are accepted against
the exact CNFs by the formally verified \texttt{cake\_lpr} checking core.\\
Semantic bridge &
Proposition~\ref{prop:encoding} proves the intended equivalence, and a separate
direct graph implementation agrees on deterministic test assignments.
The source implementation is not formally verified.\\
Catalogue &
The generation scheme is exhaustive if its implementation and nauty
canonicalisation are correct.  Its full sequence matches McKay's
maintained data, but the C generator and nauty are not proved correct in a
proof assistant.\\
System &
The compiler, operating system, file I/O, and hardware remain in the
conventional trusted computing base, except where the CakeML checker
reduces that base for the SAT-checking layer.\\
\bottomrule
\end{tabularx}
\end{center}

The replay record and this companion are bound to the following artifacts.

\begin{itemize}
  \item Current archival dependency:
  \url{https://github.com/ipitchford/z20-cochromatic}, annotated tag
  \artifact{candidate-2026-07-26}, commit
  \sha{3c7e520fdc0615f5c700761c2b1e5108dcc836e7}.
  \item Historical replay provenance: local source commit \sha{61e8300}.
  The dated receipt records that the replay originally reviewed public
  snapshot \sha{8d9829edebecdfc78384eef5b41617ddd32b11ea};
  the load-bearing source and proof objects are byte-identical at the
  archival dependency above.
  \item Replay date and platform: 26 July 2026, macOS on arm64.
  \item Replay receipt:
  \artifact{paper/z20_equals_6_replay_receipt_2026-07-26.md}.
\end{itemize}

At the pinned commit, the load-bearing public paths are:

\begin{itemize}
  \item catalogue: \path{src/gen44.c} and \path{src/gen44_result.txt};
  \item formula generator: \path{certificates/generate.py};
  \item core \(0\): \path{certificates/z20_core0.cnf},
  \path{certificates/z20_core0.drup}, and
  \path{certificates/core0.lrat};
  \item core \(1\): \path{certificates/z20_core1.cnf},
  \path{certificates/z20_core1.drup}, and
  \path{certificates/core1.lrat};
  \item package checks: \path{certificates/audit.py} and
  \path{certificates/check_drup.c}; and
  \item separate semantic/RUP replay:
  \path{verification/verify_certificate.py}.
\end{itemize}

The companion directory contains an exact path-and-hash map and a small
direct verifier for the \(19\)-vertex lower construction.  Nauty,
DRAT-trim, and \texttt{cake\_lpr} are not vendored.

\statusbox{The accurate present claim is: \emph{a complete, locally
regenerated candidate computer-assisted determination of
\(\mathrm{VR}_2(K_4)=20\), with a hand-checkable lower bound,
certificate-backed fixed-core refutations, and an explicit trusted-
computing boundary}.}

\section{Provenance and AI disclosure}

The source repository reports that the underlying mathematical and
computational work was produced by AI systems, while the human participant
selected the problem, directed the work, preserved the artifacts, and
mediated between systems.  Ian Pitchford is the repository maintainer and
publisher and performed that problem-selection and research-mediation role.
The reported division of labour is: an earlier
Hermes-driven search; subsequent instrumentation and survivor analysis;
the two-core reduction, CNF generator, DPLL solver, and DRUP certificates
from GPT-5.6 Sol; and separate verification, catalogue regeneration,
and earlier exposition by Claude Opus 5.  The present companion paper was
prepared for repository publication with OpenAI Codex from the preserved
proof manuscript, replay receipt, source code, certificates, and the sharp
lower-bound construction.  This paragraph records project provenance; it
is not an externally audited authorship determination.

\section{Licence}

To the extent copyright or related rights subsist and are held by the
repository affirmer, the original text, source, verifier, document layout,
and other original companion material are dedicated to the public domain
under CC0 1.0 Universal.  The named AI systems are provenance and credit
labels, not the legal CC0 affirmer.  Embedded font software, cited works,
linked external material, third-party tools, and model or provider
trademarks retain their own terms.  The complete legal text and scope
notice are in the repository's \artifact{LICENSE} and
\artifact{PUBLIC_DOMAIN.md}.

\appendix
\section{Exact lower-bound witness data}

With vertices ordered as
\[
  0,1,\ldots,8,9,10,\ldots,16,17,18,
\]
where \(9,17,18\) are the three independent twins, one graph6 encoding of
\(L_{19}\) is
\[
  \texttt{\detokenize{RzlXWmJpZDeJEJbDgp\EJsWnDlEJY?}}.
\]
The backslash is a literal graph6 character.  The direct verifier supplied
with the paper reconstructs the graph from the Paley rule rather than
trusting this string.  It enumerates all \(\binom{19}{4}=3{,}876\)
four-sets and confirms:
\[
\begin{array}{lr}
\text{red edges} &84,\\
\text{red }K_4\text{s} &0,\\
\text{blue }K_4\text{s} &44,\\
\text{vertex-disjoint monochromatic pair} &\text{none}.
\end{array}
\]

\small\raggedright
\begin{thebibliography}{99}

\bibitem{greenwood-gleason}
R. E. Greenwood and A. M. Gleason,
``Combinatorial relations and chromatic graphs,''
\emph{Canadian Journal of Mathematics} \textbf{7} (1955), 1--7.
\url{https://doi.org/10.4153/CJM-1955-001-4}.

\bibitem{jent-osborn-zhang}
E. Jent, S. Osborn, and P. Zhang,
``Extending Ramsey numbers for connected graphs of size 3,''
\emph{Symmetry} \textbf{16} (2024), Article 1092.
\url{https://doi.org/10.3390/sym16081092}.

\bibitem{jent-zhang}
E. Jent and P. Zhang,
``Multiple Ramsey numbers of disconnected graphs of size 3,''
\emph{Journal of Combinatorial Mathematics and Combinatorial Computing}
\textbf{129} (2026), 17--31.
\url{https://doi.org/10.61091/jcmcc129-02}.

\bibitem{mckay-catalogue}
B. D. McKay,
``Ramsey graphs,'' Combinatorial Data collection,
\url{https://users.cecs.anu.edu.au/~bdm/data/ramsey.html},
accessed 26 July 2026.

\bibitem{mckay-piperno}
B. D. McKay and A. Piperno,
``Practical graph isomorphism, II,''
\emph{Journal of Symbolic Computation} \textbf{60} (2014), 94--112.
\url{https://doi.org/10.1016/j.jsc.2013.09.003}.

\bibitem{paley}
R. E. A. C. Paley,
``On orthogonal matrices,''
\emph{Journal of Mathematics and Physics} \textbf{12} (1933), 311--320.
\url{https://doi.org/10.1002/sapm1933121311}.

\bibitem{drat-trim}
N. Wetzler, M. J. H. Heule, and W. A. Hunt, Jr.,
``DRAT-trim: Efficient checking and trimming using expressive clausal
proofs,'' in \emph{Theory and Applications of Satisfiability Testing --
SAT 2014}, LNCS 8561, 422--429.
\url{https://doi.org/10.1007/978-3-319-09284-3_31}.

\bibitem{cake-lpr}
Y. K. Tan, M. J. H. Heule, and M. O. Myreen,
``Verified propagation redundancy and compositional UNSAT checking in
CakeML,''
\emph{International Journal on Software Tools for Technology Transfer}
\textbf{25} (2023), 167--184.
\url{https://doi.org/10.1007/s10009-022-00690-y}.

\bibitem{z20repo}
I. Pitchford,
\emph{\(z(20)=6\): a candidate computer-assisted proof},
public artifact repository, annotated tag
\artifact{candidate-2026-07-26}, commit
\sha{3c7e520fdc0615f5c700761c2b1e5108dcc836e7},
\url{https://github.com/ipitchford/z20-cochromatic},
accessed 26 July 2026.

\end{thebibliography}
\normalsize

\end{document}
